\section{Variation at equality and localization of the area term}
\label{q4:var:section}

Throughout this section, $(\Omega,g,K)$ satisfies the hypotheses of
Definition~\ref{q4:def:horizon}, and equality holds in
Theorem~\ref{q4:thm:numerical-provider}. Thus $\Omega$ has one end and its only boundary
$S$ is a connected, outermost, outer area-minimizing future MOTS. We use the
full enclosing-cut convention of Definition~\ref{q4:def:cuts}. Set
\begin{equation}
\begin{gathered}
 A=\Area_g(S),\qquad m=\sqrt{E^2-|P|^2},\qquad
 r_0=(A/\omega)^{1/3}=\sqrt{2m},\\
 b^0=E/m,\qquad b^i=-P_i/m,\qquad
 \mathfrak m(a)=\tfrac12(a/\omega)^{2/3},\qquad
 c=6\omega\mathfrak m'(A)=2/r_0.
\end{gathered}
\label{q4:var:constants}
\end{equation}
The strict future timelikeness needed here is already part of
Theorem~\ref{q4:thm:numerical-provider}. For nearby data with charges
$(\widetilde E,\widetilde P)$, define the fixed linear charge functional
\[
 \mathcal E(\widetilde g,\widetilde K)
   =b^0\widetilde E+\sum_{i=1}^4b^i\widetilde P_i.
\]
The reverse Cauchy--Schwarz inequality on the future timelike cone gives
\begin{equation}
 \mathcal E(\widetilde g,\widetilde K)
 \ge \sqrt{\widetilde E^2-|\widetilde P|^2}.
\label{q4:var:supporting-charge}
\end{equation}
Equality holds at the original data. In particular, whenever the varied
data satisfy the numerical theorem, their enclosing infimum
$a(\widetilde g)$ satisfies
\begin{equation}
 \mathcal E(\widetilde g,\widetilde K)
       -\mathfrak m(a(\widetilde g))\ge0.
\label{q4:var:nonnegative-defect}
\end{equation}
The varied data need not preserve outermostness or outer area minimization:
those assumptions are absent from the numerical theorem.

The right-derivative formula in
Proposition~\ref{q4:enc:area-derivative} accounts for all minimizing
frontiers without choosing a differentiable family. It gives the
positive-measure identity of Proposition~\ref{q4:var:multiplier}.
Compact normal tests and outermostness then localize its area term to $S$
in Proposition~\ref{q4:var:area-localization}, supplying the identity used
to construct the causal adjoint field in Theorem~\ref{q4:adj:causal-field}.

\subsection{Constraint derivatives and a strict direction}

On the closed unit ball bundle of $g$, introduce
\begin{equation}
\begin{gathered}
 C(x,v)=2(\mu(x)+J_x(v))
       =R_g+\tau^2-|K|_g^2+2J_x(v),\qquad |v|_g\le1,\\
 \mathcal A=\{(x,v):C(x,v)=0\}.
\end{gathered}
\label{q4:var:active-set}
\end{equation}
DEC is equivalent to $C\ge0$ on this bundle. When the metric varies,
we identify its unit ball with the original one by the positive symmetric
inverse square root. More precisely, if
$g_s(V,W)=g(G_sV,W)$, then $I_s=G_s^{-1/2}$ sends the $g$ unit ball
isometrically to the $g_s$ unit ball. For a variation
$H=(h,p)=(\dot g_0,\dot K_0)$, all derivatives $C'_H$ below include this
identification; in particular,
\[
 \dot I_0v=-\tfrac12h^\sharp v.
\]

\begin{lemma}[Conformal constraint identities]
\label{q4:var:conformal}
For a positive smooth function $f$, let $g_f=f^2g$ and $K_f=fK$.
Using $f^{-1}v$ in the varied unit ball, one has
\begin{align}
 f^2 C_f(x,f^{-1}v)
   &=C(x,v)+6f^{-1}\bigl[-\Delta_g f+K(\nabla f,v)\bigr],
       \label{q4:var:conformal-C}\\
 f\theta_{+,f}
   &=\theta_++3\partial_\nu\log f.
       \label{q4:var:conformal-theta}
\end{align}
Here $\nu$ points into the exterior and $\Delta_g=\operatorname{div}_g\nabla$.
\end{lemma}

\begin{proof}
Put $\varphi=\log f$. The scalar-curvature and tensor transformations in
dimension four are
\[
 R_{f^2g}=f^{-2}(R_g-6\Delta_g\varphi-6|d\varphi|_g^2),
 \quad \tau_f=f^{-1}\tau,
 \quad |K_f|_{f^2g}^2=f^{-2}|K|_g^2.
\]
Since $\Delta\varphi+|d\varphi|^2=f^{-1}\Delta f$, they give
$f^2\mu_f=\mu-3f^{-1}\Delta_gf$. For the momentum, set
$\pi=K-\tau g$, so $\pi_f=f\pi$. The connection difference is
\[
 \widehat\Gamma^a_{ij}-\Gamma^a_{ij}
 =\delta_i^a\varphi_j+\delta_j^a\varphi_i-g_{ij}\varphi^a.
\]
Contracting the covariant derivative of $f\pi$ yields
\[
 fJ_{f,i}=J_i+3\pi_{ij}\varphi^j-(\tr_g\pi)\varphi_i
          =J_i+3K_{ij}\varphi^j,
\]
because $\tr_g\pi=-3\tau$. In particular every trace-gradient term
cancels; no extra asymptotic condition on $\tau$ is used. Evaluating on
the identified vector $f^{-1}v$ gives
\[
 f^2J_f(f^{-1}v)=J(v)+3f^{-1}K(\nabla f,v).
\]
Adding twice the density and momentum identities proves
\eqref{q4:var:conformal-C}. Finally
$\nu_f=f^{-1}\nu$,
$H_{f^2g}=f^{-1}(H_g+3\partial_\nu\log f)$, and
$\tr_{TS,f^2g}(fK)=f^{-1}\tr_{TS,g}K$.
Their sum proves \eqref{q4:var:conformal-theta} with the same oriented
normal.
\end{proof}

\begin{lemma}[Strict conformal direction]
\label{q4:var:strict-direction}
Choose
\[
 1<q_0<\min(q,2),\qquad 0<\delta<\min(q_0,1),
 \qquad w=r^{-4-\delta},
\]
where the end radius is extended to a positive smooth function on $\Omega$.
There is a smooth nonnegative function $\phi$, smooth up to $S$, such that
\begin{equation}
 \partial_\nu\phi=-1\quad\hbox{on }S,\qquad
 -\Delta_g\phi-|K|_g|\dd\phi|_g\ge w,\qquad
 \phi=O_2(r^{-2}),\qquad -\Delta_g\phi/w\longrightarrow1.
\label{q4:var:phi-properties}
\end{equation}
Its limiting Euclidean sphere gradient flux exists and is finite.
For
\[
 Q=(2\phi g,\phi K)
\]
one consequently has
\begin{equation}
 C'_Q\ge6w\quad\hbox{on }\mathcal A,
 \qquad -\theta'_Q=3\quad\hbox{on }S.
\label{q4:var:strict-derivatives}
\end{equation}
Moreover, uniformly for all $|v|_g\le1$ on the end,
$C'_Q/w\longrightarrow6$.
\end{lemma}

\begin{proof}
Choose a smooth majorant $d_0\ge |K|_g$ which on the end is a sufficiently
large radial multiple of $r^{-1-q_0}$. It can be chosen with all symbol
bounds there. We solve
\begin{equation}
 -\Delta_g\phi
   =d_0\sqrt{|\dd\phi|_g^2+r^{-6}}+w,
 \qquad \partial_\nu\phi=-1\text{ on }S,
 \qquad \phi\longrightarrow0\text{ at infinity}.
\label{q4:var:phi-equation}
\end{equation}

First truncate at a large coordinate sphere $S_R$, imposing $\phi=0$ on
$S_R$, and replace $d_0$ in the equation by $t d_0$, $0\le t\le1$.
The inner and outer boundary components are disjoint smooth hypersurfaces;
there is no corner where the two boundary conditions meet. The $t=0$
problem is the invertible mixed Laplace problem. Every linearization in
$\phi$ adds to $-\Delta_g$ a smooth drift of norm at most $d_0$, with
homogeneous Neumann data at $S$ and homogeneous Dirichlet data at $S_R$.
The maximum and boundary point principles give uniqueness, and the mixed
elliptic Fredholm theory then gives invertibility. These are the usual
local elliptic and boundary estimates on a fixed smooth truncated domain;
see \cite{GilbargTrudinger2001}.

Solutions are nonnegative. Indeed a negative minimum cannot occur in the
interior because the right side of the equation is positive. At $S$, the
outward normal of the truncated domain is $-\nu$, so the prescribed
outward derivative is $+1$, also excluding such a minimum by the boundary
point principle. The outer Dirichlet value is zero.

For completeness, the continuation bounds do not require a monotonicity
assumption in the unknown. Fix $p>4$. The mixed $W^{2,p}$ estimate, the
linear growth of the right side in $\dd\phi$, and first-derivative
interpolation give, on this fixed truncation,
\[
 \|\phi\|_{W^{2,p}}\le C\bigl(1+\|\phi\|_{L^\infty}\bigr),
\]
uniformly in $t$. If the supremum were unbounded, division by that
supremum and compactness in $C^1$ would produce a nonnegative function
$z$ with supremum one satisfying
\[
 -\Delta_g z=t_*d_0|\dd z|_g,
 \qquad \partial_\nu z=0\text{ on }S,
 \qquad z=0\text{ on }S_R
\]
for some $t_*\in[0,1]$. At points where $\dd z\ne0$ the right side is a
bounded drift applied to $\dd z$, and it is zero otherwise. The strong
maximum and boundary point principles for this bounded-drift equation
contradict the positive maximum. The resulting $W^{2,p}$ and Schauder
bounds close the continuation argument and give a smooth solution at
$t=1$ on every sufficiently large truncation.

We next make the bounds independent of $R$. For
$W=r^{-2}-r^{-2-\delta}$, the Euclidean identity
\[
 -\Delta_\delta W=(2+\delta)\delta r^{-4-\delta}
\]
and the AF estimates imply, beyond a fixed radius $R_1$,
\begin{equation}
 -\Delta_g W-d_0|\dd W|_g\ge c_1w>0.
\label{q4:var:tail-barrier}
\end{equation}
The metric errors and the drift term are smaller than $w$ because
$\delta<q_0<q$. Also $d_0r^{-3}=o(w)$.
If $M_R$ is the supremum of the truncated solution on the fixed core
bounded by $S_{R_1}$, a sufficiently large fixed multiple of
$(1+M_R)W$ is a supersolution on $R_1\le r\le R$. To see this, use
$\sqrt{a^2+b^2}\le a+b$ and \eqref{q4:var:tail-barrier}; choose the multiple
also to dominate $M_R$ on $S_{R_1}$. It dominates the zero outer value.
Comparison is valid by writing the difference of the two gradient
nonlinearities as a bounded drift. Therefore
\begin{equation}
 0\le\phi_R\le C(1+M_R)W\qquad (R_1\le r\le R).
\label{q4:var:tail-bound}
\end{equation}

If $M_R$ diverged along a sequence, normalize by $M_R$. Local mixed
estimates near $S$ and interior estimates elsewhere give a subsequential
limit $z$ on every fixed compact set. More precisely,
\eqref{q4:var:tail-bound} bounds $\phi_R/M_R$ on a neighborhood of the fixed
closed core. Local Neumann $W^{2,p}$ estimates at $S$, and interior
$W^{2,p}$ estimates away from it, give uniform bounds there for $p>4$.
Compactness in $C^1$ on the closed core retains both
$\sup_{\mathrm{core}}z=1$ and $\partial_\nu z=0$.
The limit satisfies
$-\Delta_gz=d_0|\dd z|_g$, homogeneous Neumann data on $S$, and,
by \eqref{q4:var:tail-bound}, $z\to0$ at infinity. Its positive global maximum
is attained in a compact set. The same strong maximum and boundary point
principles again give a contradiction. Thus $M_R$ is bounded, and
exhaustion produces a smooth solution of \eqref{q4:var:phi-equation} with
$\phi=O(r^{-2})$.

Here is the derivative control needed later, using only the stated
derivatives of the original data. On a fixed annulus in the $y$ variables
put $g_\rho(y)=(g_{ij}(\rho y))$, $U_\rho(y)=\rho^2\phi(\rho y)$, and
$a_\rho(y)=\rho d_0(\rho y)$. The exact rescaled equation is
\begin{equation}
 -\Delta_{g_\rho}U_\rho
 =a_\rho\sqrt{|\dd U_\rho|_{g_\rho}^2+|y|^{-6}}
       +\rho^{-\delta}|y|^{-4-\delta}.
\label{q4:var:rescaled-phi}
\end{equation}
The metric components have uniform $C^{1,1}$ bounds from $g-\delta=O_2$.
Consequently both the principal coefficients of this Laplacian and its
first-order coefficients have uniform $C^{0,\alpha}$ bounds for every
fixed $\alpha<1$. No third derivative of $g$ is needed. The chosen
radial $d_0$ gives $a_\rho=O(\rho^{-q_0})$ with uniform symbol bounds,
independently of derivatives of $K$. First-derivative interpolation in
the local $W^{2,p}$ estimates gives a uniform $W^{2,p}$ bound on a smaller
annulus. Taking $p>4$ controls $\dd U_\rho$ in $C^{0,\alpha}$ for some
$\alpha>0$. The positive $|y|^{-6}$ term is bounded away from zero there,
so the entire right side of \eqref{q4:var:rescaled-phi} is uniformly
$C^{0,\alpha}$. The usual Schauder estimate, including the first-order
Laplacian coefficients with these $C^{0,\alpha}$ bounds, now gives a
uniform $C^{2,\alpha}$ bound on a further smaller annulus. Scaling back
proves $\phi=O_2(r^{-2})$ without an additional metric or tensor
falloff derivative.
The equation now yields
\[
 d_0\sqrt{|\dd\phi|_g^2+r^{-6}}=O(r^{-4-q_0})=o(w),
\]
which proves both the differential inequality and the last limit in
\eqref{q4:var:phi-properties}. The right side of \eqref{q4:var:phi-equation} is
integrable. The divergence theorem consequently gives a finite limiting
physical sphere flux. Its difference from the Euclidean sphere flux tends
to zero, since $g-\delta=O(r^{-q})$ and $\dd\phi=O(r^{-3})$.

Finally, differentiating Lemma~\ref{q4:var:conformal} at $f=1+s\phi$ gives
\[
 C'_Q=-2\phi C+6[-\Delta_g\phi+K(\nabla\phi,v)],
 \qquad \theta'_Q=-\phi\theta_++3\partial_\nu\phi.
\]
On $\mathcal A$ the first term is zero, and on $S$ we have
$\theta_+=0$. This proves \eqref{q4:var:strict-derivatives}.
On the full end ball bundle, $C=O(r^{-2-q})$,
$\phi C=O(r^{-4-q})=o(w)$, and
$K(\nabla\phi,v)=O(r^{-4-q})=o(w)$. Thus the normalized derivative tends
uniformly to six.
\end{proof}

\subsection{Feasible variations and a measure identity}

\begin{definition}[Variation space]
\label{q4:var:variation-space}
Let $\mathcal V_0$ be the real vector space generated by all smooth compactly
supported pairs $(h,p)$ up to $S$ and the following five pairs, cut off
smoothly to the end. The scalar pair is
\[
 h_{ij}=r^{-2}\delta_{ij},\qquad p=0.
\]
For each vector $e$ in a fixed basis of $\R^4$, the momentum pair has
$h=0$ and is defined by
\begin{equation}
 p-(\tr_\delta p)\delta
 =r^{-3}\bigl(e\otimes\widehat x+\widehat x\otimes e
                  -(e\cdot\widehat x)\delta\bigr),
 \qquad \widehat x=x/r.
\label{q4:var:momentum-prototype}
\end{equation}
Set $\mathcal V=\R Q+\mathcal V_0$. We write its elements as
$H=aQ+H_0$, where $H_0=(h_0,p_0)\in\mathcal V_0$.
\end{definition}

\begin{lemma}[End behavior of the variation space]
\label{q4:var:end-variations}
The five prototypes have independent charge derivatives. For every fixed
$H_0\in\mathcal V_0$,
\begin{equation}
 C'_{H_0}=O(r^{-4-q})=o(w)
\label{q4:var:prototype-constraint-decay}
\end{equation}
uniformly on the full end ball bundle. The coefficient $a$ in
Definition~\ref{q4:var:variation-space} is uniquely determined, and
\begin{equation}
 C'_H/w\longrightarrow6a
\label{q4:var:end-value}
\end{equation}
uniformly there. All elements of $\mathcal V$ have well-defined charge
derivatives.
\end{lemma}

\begin{proof}
Trace reversal on symmetric tensors in dimension four is invertible. The
right side of \eqref{q4:var:momentum-prototype}, written as
\[
 D_{ij}=r^{-4}\bigl(e_i x_j+x_i e_j-(e\cdot x)\delta_{ij}\bigr),
\]
satisfies $\partial_jD_{ij}=0$ and $D_{ij}\widehat x^j=r^{-3}e_i$.
The scalar prototype satisfies the Euclidean linearized scalar constraint
because $\Delta_\delta r^{-2}=0$ off the origin. With the ADM normalizations
\eqref{q4:intro:energy}--\eqref{q4:intro:momentum}, the scalar prototype
has $E'=1$, $P'=0$, and the momentum
prototype for $e$ has $E'=0$, $P'=e/3$. Background corrections to these
fluxes vanish by the AF estimates, so the derivatives are independent.

On the end, $h_0=O_2(r^{-2})$ and $p_0=O_1(r^{-3})$. Their Euclidean
linearized scalar and momentum constraints vanish. Each remaining linear
constraint term has a factor from the decaying background: for example,
$(g-\delta)\partial^2h_0$, $\partial g\,\partial h_0$,
$\partial^2g\,h_0$, $Kp_0$, or $(g-\delta)\partial p_0$ has order at most
$r^{-4-q}$. The variation of the vector identification contributes
$-J(h_0^\sharp v)$, of the same order. This proves
\eqref{q4:var:prototype-constraint-decay}. Lemma~\ref{q4:var:strict-direction}
then proves \eqref{q4:var:end-value} and the uniqueness of $a$: an element of
$\mathcal V_0$ cannot have the normalized end derivative of $Q$.

Compact variations have zero charge derivatives. The prototypes have the
fluxes just computed. For $Q$, the momentum derivative is zero and
\[
 E'_Q=-\frac1\omega\lim_{R\to\infty}
              \int_{|x|=R}\partial_r\phi\,\dd A_\delta,
\]
which is finite by Lemma~\ref{q4:var:strict-direction}. Products with
background errors have zero limiting flux. Linearity completes the proof.
\end{proof}

Let $\mathcal T$ be the compact space of minimizing frontiers for $g$,
with their filled sets, from Proposition~\ref{q4:enc:area-derivative}.
For $T\in\mathcal T$, write
\begin{equation}
 a'_T(h)=\tfrac12\int_T\tr_T h\,\dd A_g.
\label{q4:var:fixed-area-derivative}
\end{equation}
This is a continuous function of $T$ for each fixed variation $h$.

\begin{lemma}[Strict linear inequalities give feasible nonlinear data]
\label{q4:var:feasible-path}
Suppose $H=aQ+H_0\in\mathcal V$ has $a>0$ and, for some $\varepsilon>0$,
\[
 C'_H\ge\varepsilon w\text{ on }\mathcal A,\qquad
 -\theta'_H>0\text{ on }S.
\]
Then, for every sufficiently small $s>0$,
\begin{equation}
 g_s=e^{2as\phi}(g+sh_0),\qquad
 K_s=e^{as\phi}(K+sp_0)
\label{q4:var:nonlinear-path}
\end{equation}
satisfies every hypothesis of Theorem~\ref{q4:thm:numerical-provider}. Its boundary
has strictly negative future expansion. The path has derivative $H$, is
uniformly comparable to $g$, has uniform large-sphere barriers, and its
charges are differentiable at $s=0$.
\end{lemma}

\begin{proof}
Positivity of the metric and uniform comparability hold for small $s$,
since $h_0$ and $\phi$ are bounded relative to $g$. They also give
completeness with the boundary included. The topology, smoothness, end
count and orientability remain those of $\Omega$.

We check DEC on the whole ball of test vectors. First use the intermediate
metric and tensor
$\overline g_s=g+sh_0$, $\overline K_s=K+sp_0$, identifying its vector
ball with that of $g$ by $\overline I_s$. For sufficiently large $r$,
Taylor expansion and the Euclidean cancellations in
Lemma~\ref{q4:var:end-variations} give, uniformly in $|v|_g\le1$,
\[
 \overline C_s(x,\overline I_sv)
    =C(x,v)+O(sr^{-4-q})+O(s^2r^{-6}).
\]
Here the remainder is uniform in $r$, $s$ in a fixed small interval, and
the whole vector ball. To check all its components, the coordinate
constraint expressions have the schematic forms
\begin{align*}
 R_g&=g^{-1}\partial^2g+\operatorname{smooth}(g^{-1})(\partial g)^2,\\
 J&=\operatorname{smooth}(g^{-1})\partial K
       +\operatorname{smooth}(g^{-1})(\partial g)K,
\end{align*}
and $2\mu-R_g$ is quadratic in $K$, with smooth coefficients in $g^{-1}$.
Each second $s$ derivative on the intermediate path is $O(r^{-6})$:
the largest terms are $h_0\partial^2h_0$, $(\partial h_0)^2$, $p_0^2$,
$h_0\partial p_0$, and $(\partial h_0)p_0$. Terms containing an original
$K$, $\partial K$, or background metric derivative decay faster. This
calculation includes all occurrences of $\tau=g^{ij}K_{ij}$, without
a separate trace estimate. The vector identification also has the uniform
matrix expansion
\[
 \overline I_s=(\Id+s h_0^\sharp)^{-1/2}
       =\Id-\tfrac{s}{2}h_0^\sharp+O(s^2r^{-4});
\]
the eigenvalues lie in a fixed positive interval, so the matrix-function
Taylor bound is uniform. Multiplication by the original or varied
momentum covector keeps the same remainder order. In fact the preceding
calculation gives the componentwise estimates
\begin{equation}
\begin{aligned}
 \overline\mu_s&=\mu+O(sr^{-4-q})+O(s^2r^{-6}),\\
 \overline J_s&=J+O(sr^{-4-q})+O(s^2r^{-6}).
\end{aligned}
\label{q4:var:component-remainders}
\end{equation}

Apply Lemma~\ref{q4:var:conformal} to these intermediate data with
$f_s=e^{as\phi}$. After multiplication by the positive factor $f_s^2$, the
constraint for the final data is
\begin{equation}
 C(x,v)+s\bigl(6aw+o(w)\bigr)+O(s^2r^{-6}),
\label{q4:var:end-feasibility-expansion}
\end{equation}
again uniformly for the full vector ball. Indeed,
$-\Delta_g\phi/w\to1$ and $K(\nabla\phi,v)=o(w)$. Changing the metric,
tensor, and vector identification in these conformal differential terms
adds only the above decay orders, while the extra exponential term uses
$|\dd\phi|^2=O(r^{-6})$. Since $a>0$ and
$r^{-6}=O(w)$, one can first choose a fixed large radius so the coefficient
of $s$ in \eqref{q4:var:end-feasibility-expansion} is positive, and then
choose $s$ small so its remainder is absorbed. DEC follows on this fixed
far end, using $C\ge0$.

The remaining base region and its closed unit ball bundle are compact.
The derivative of the constraint is strictly positive on its original
zero set. By continuity it is positive on a neighborhood of that zero
set; on the complement the original constraint has a positive minimum.
A uniform Taylor expansion therefore gives $C_s\ge0$ everywhere in the
compact region for small positive $s$. The same compactness argument on
$S$, where $\theta_+=0$, gives $\theta_{+,s}<0$.

The new densities are integrable componentwise, not merely after
contraction with active vectors. The exact conformal formulas give
\begin{align*}
 f_sJ_s&=\overline J_s+
            3as\overline K_s(\nabla_{\overline g_s}\phi,\cdot),\\
 f_s^2\mu_s&=\overline\mu_s-3as\Delta_{\overline g_s}\phi
                          -3a^2s^2|\dd\phi|_{\overline g_s}^2.
\end{align*}
The extra covector in the first line is
$O(sr^{-4-q})+O(s^2r^{-6})$. Also
$\Delta_{\overline g_s}\phi=\Delta_g\phi+O(sr^{-6})$, with
$\Delta_g\phi=O(w)$. Together with
\eqref{q4:var:component-remainders}, these identities express the new
densities as bounded smooth scaling factors times the original integrable
densities, plus integrable errors of orders $r^{-4-\delta}$,
$r^{-4-q}$, and $r^{-6}$. Uniform metric comparability controls both norms
and volume elements. The differentiated falloff has
exponent $\min(q,2)>1$, uniformly for $s$ in a small fixed interval, so
large coordinate spheres remain mean convex uniformly.

The charge limits can also be checked before differentiation. Write
$\pi=K-\tau g$. Since conformal trace reversal gives
$\pi_s=f_s(\overline K_s-\overline\tau_s\overline g_s)$, expansion on the
end yields
\[
 \pi_s=\pi+s\bigl(p_0-(\tr_\delta p_0)\delta\bigr)
          +O(sr^{-3-q})+O(s^2r^{-5}).
\]
The terms with the original trace are of size
$h_0K=O(r^{-3-q})$; multiplication by $f_s-1=O(sr^{-2})$ adds only the
displayed error orders. All errors have zero limiting sphere flux.
Therefore $P_s=P+sP'_{H_0}$ in the prescribed coordinates, without an
extra decay assumption on $\tau$. Similarly,
\[
 g_s-g-s(h_0+2a\phi g)=O_1(s^2r^{-4}),
\]
whose energy flux vanishes. The finite flux of $\phi$ and the scalar
prototype thus give $E_s=E+sE'_H$. In particular both charge functions
exist and are differentiable with the derivatives computed in
Lemma~\ref{q4:var:end-variations}. The path and its
first two parameter derivatives are uniformly bounded relative to $g$,
as required by Proposition~\ref{q4:enc:area-derivative}. This verifies all
claims and all hypotheses of the numerical theorem.
\end{proof}

\begin{proposition}[Multiplier identity at equality]
\label{q4:var:multiplier}
There exist a probability measure $\varpi$ on $\mathcal T$, a nonnegative
finite measure $\eta$ on $S$, a nonnegative locally finite measure
$\Lambda$ on $\mathcal A$, and $z\ge0$, with
$\int_{\mathcal A}w\,\dd\Lambda<\infty$, such that for every
$H=aQ+H_0=(h,p)\in\mathcal V$,
\begin{equation}
 6\omega\mathcal E'(H)-c\int_{\mathcal T}a'_T(h)\,\dd\varpi(T)
   =\langle C'_H,\Lambda\rangle
      -\int_S\theta'_H\,\dd\eta+za.
\label{q4:var:measure-identity}
\end{equation}
The constraint pairing is absolutely defined, since $C'_H/w$ is bounded
on the compactified active set.
\end{proposition}

\begin{proof}
Compactify $\mathcal A$ by one end point. If $\mathcal A$ is already
compact, adjoin an isolated point; if it is empty, use just that point.
The ball fibers are compact and the base has one end, so this is a compact
Hausdorff space, denoted $\overline{\mathcal A}$. By
Lemma~\ref{q4:var:end-variations}, $C'_H/w$ extends continuously to it with
value $6a$ at the added point. On the disjoint compact union
\[
 \mathcal K=\overline{\mathcal A}\sqcup S\sqcup\mathcal T
\]
consider the linear map to $C(\mathcal K)$ whose three components are
\begin{equation}
 H\longmapsto
 \left(\frac{C'_H}{w},\ -\theta'_H,\
                   c a'_T(h)-6\omega\mathcal E'(H)\right).
\label{q4:var:separation-map}
\end{equation}

Its image does not meet the open cone of functions strictly positive
everywhere. Otherwise its value at the added point gives $a>0$, and
Lemma~\ref{q4:var:feasible-path} supplies actual feasible data
\eqref{q4:var:nonlinear-path}. Let $a(s)$ be their enclosing infimum. The
area derivative formula in Proposition~\ref{q4:enc:area-derivative} gives
\[
 a'(0+)=\min_{T\in\mathcal T}a'_T(h).
\]
Strict positivity of the third component in \eqref{q4:var:separation-map},
and compactness of $\mathcal T$, imply
\[
 \left.\frac{\dd}{\dd s}\right|_{0+}
   6\omega\bigl(\mathcal E(g_s,K_s)-\mathfrak m(a(s))\bigr)
   =6\omega\mathcal E'(H)-c\min_{T\in\mathcal T}a'_T(h)<0.
\]
The expression starts at zero and is nonnegative by
\eqref{q4:var:nonnegative-defect}, a contradiction.

Hahn--Banach separation of a linear subspace from this open convex cone
gives a nonzero positive continuous functional on $C(\mathcal K)$
annihilating the image. The subspace need not be closed: its closure is
still disjoint from the open cone. The Riesz representation theorem
identifies the functional with a nonzero finite positive measure on
$\mathcal K$.

Its mass on $\mathcal T$ is positive. If that mass were zero, testing $Q$
would give a strictly positive value on every remaining component:
$C'_Q/w\ge6$ on $\mathcal A$, the added-point value is six, and
$-\theta'_Q=3$ on $S$. This would contradict annihilation. Normalize the
mass on $\mathcal T$ to one and call the resulting measure there
$\varpi$. Call its restriction to $S$ $\eta$. On $\mathcal A$, divide
the remaining measure by the positive function $w$ to obtain $\Lambda$;
this is locally finite and has finite weighted mass. If the normalized
added-point mass is $\lambda_\infty$, put $z=6\lambda_\infty$.
Rearranging annihilation of \eqref{q4:var:separation-map} is exactly
\eqref{q4:var:measure-identity}.
\end{proof}

The separation argument has supplied a probability measure on all
minimizing enclosures, not a fixed-horizon first variation. Our next task
is to remove its interior support. We do this before extracting a
homogeneous adjoint field: compact normal motions test the trace of $K$
on the minimizing frontiers, and the resulting MOTS equation brings
outermostness into the argument.

\subsection{Normal graph tests}

The measure $\varpi$ may initially be supported on several enclosures of
the same area. The next argument uses variations of the original data to
show that it is supported on the original horizon. It is local in the open
exterior and assumes no complete ambient development.

\begin{lemma}[Compact normal tests annihilating active constraints]
\label{q4:var:normal-tests}
For every $\ell\in C^\infty_c(\operatorname{int}\Omega)$, there are smooth
variations $H_j=(h_j,p_j)\in\mathcal V_0$, all supported in one fixed
compact subset of the open exterior, such that
\[
 h_j=2\ell K,\qquad C'_{H_j}\longrightarrow0
\]
uniformly on the active rays over that compact set.
\end{lemma}

\begin{proof}
Choose a relatively compact open neighborhood of
$\operatorname{supp}\ell$ in $\operatorname{int}\Omega$ and apply
Lemma~\ref{lem:variation:dust-tests} there in spatial dimension four,
with $s=\ell$ and $\delta=1/j$. The present densities already satisfy
$2\mu=R_g+(\tr_gK)^2-|K|_g^2$ and
$J=\operatorname{div}_g(K-(\tr_gK)g)$, with $\mu\ge|J|_g$.
The smoothness and future-normal convention
\eqref{q4:eq:intro-second-fundamental-form} are precisely those of the
lemma. Equation~\eqref{q4:var:active-set} defines the same constraint
$C=2(\mu+J(v))$ and the same isometric ray transport, so its derivative
is the derivative $\mathfrak C'$ in that lemma.

The resulting pairs have metric component $2\ell K$ and fixed compact
support in the open exterior. They therefore belong to $\mathcal V_0$
by Definition~\ref{q4:var:variation-space}, and the lemma gives the
asserted uniform convergence, including active rays at which $\mu=0$.
Their charge, boundary, and end-coefficient variations vanish because
the support stays strictly inside the open exterior.
\end{proof}

\begin{proposition}[Localization of the area multiplier]
\label{q4:var:area-localization}
For $\varpi$-almost every element of $\mathcal T$, its frontier equals $S$,
where $\varpi$ is the measure in Proposition~\ref{q4:var:multiplier}.
Consequently, for every variation $H=(h,p)$ in $\mathcal V$,
\begin{equation}
 \int_{\mathcal T}a'_T(h)\,\dd\varpi(T)
       =a'_S(h)=\tfrac12\int_S\tr_S h\,\dd A_g.
\label{q4:var:localized-area}
\end{equation}
In particular the multiplier identity becomes
\begin{equation}
 6\omega\mathcal E'(H)-c\,a'_S(h)
  =\langle C'_H,\Lambda\rangle
       -\int_S\theta'_H\,\dd\eta+za.
\label{q4:var:localized-identity}
\end{equation}
\end{proposition}

\begin{proof}
Apply Equation~\eqref{q4:var:measure-identity} to the variations in
Lemma~\ref{q4:var:normal-tests}. Their charge, boundary, and end-coefficient
terms vanish, and $\Lambda$ is finite over their fixed compact support.
Since the metric component is $2\ell K$ and $c>0$, uniform convergence
of the active constraint derivatives gives
\begin{equation}
 \int_{\mathcal T}\int_T\ell\,\tr_T K\,\dd A_g\,\dd\varpi(T)=0
 \qquad\text{for every }\ell\in C^\infty_c(\operatorname{int}\Omega).
\label{q4:var:trace-test}
\end{equation}

Propositions~\ref{q4:enc:minimizing-hull} and
\ref{q4:enc:common-plane} give exactly the geometry needed for
Proposition~\ref{prop:variation:smooth-localization}: the full frontiers
have area $A$, lie in one compact set, have smooth minimal free parts
and $C^{1,1}$ contact graphs, and carry a continuous common tangent
plane away from $S$. The surface-area kernel is measurable by the
tangent-plane measure topology. Their end sides are connected, and
each entire smooth frontier is an admissible full cut; no component or
contact area is discarded.

Definition~\ref{q4:def:horizon} supplies connectedness and marginality of
$S$, outer area-minimality, and the wholly-interior outermostness
condition in case \textup{(a)} of
Proposition~\ref{prop:variation:smooth-localization}. Hence almost every
frontier equals $S$. Its area-derivative conclusion gives
Equation~\eqref{q4:var:localized-area}; substituting this into
Equation~\eqref{q4:var:measure-identity} gives
Equation~\eqref{q4:var:localized-identity}.
\end{proof}

